% \section{Introduction}
% Generative models have recently met or surpassed human-level standards on benchmarks across a range of tasks \citep{nori2023capabilitiesgpt4medicalchallenge, katz2024gpt, dubey2024llama}. 
% While these claims warrant skepticism and further robustness evaluation \citep{ness2024medfuzzexploringrobustnesslarge}, the impressive capabilities have created a competitive environment for training state-of-the-art models and have inspired the development of complementary methods to adapt models to particular use cases.
% For example, quantizing the weights of a neural model enables a trade-off of model precision with vRAM and disk space \citep{gholami2022survey} while methods such as Low Rank Adaptation (LoRA) \citep{hu2021lora} and prompt-tuning \citep{lester2021power} enable compute- and data-efficient model adaptation. 
% Other methods such as retrieval-augmented generation \citep{rag}, model merging \citep{model-merging}, constrained decoding \citep{constrained-decoding}, etc. \citep{dettmers2024qlora, graphrag}, have similarly contributed to the rapid development of a large population of diverse and accessible models.

% Each model in the population has an accompanying set of covariates -- scores on benchmarks, training mixture proportions, model safety scores, probability of hallucination, etc. -- that are a function of the model, the training set, the architecture, the retrieval database, or derivatives thereof.
% For a given model, the covariate of interest or the function to calculate it may not be available to the user.
% For example, a proprietary model may have been trained on copyrighted data and the indicator for whether or not the model has had access to the data is unknown.
% Methods for predicting model-level covariates are necessary in these settings, and others, to fully understand the behavior and properties of a model.

% In this paper we extend recent theoretical results for embedding-based representations of generative models \citep{aranyak} to statistical inference settings.
% In particular, our results show that the embedding-based representations of a collection of models can be used for consistent inference for a wide class of inference problems.
% We demonstrate the effectiveness of the representations for three downstream inference tasks, including predicting the presence of sensitive information in a model's training mixture and predicting a model's safety. We include empirical investigations of  performance sensitivity to hyperparameters required to induce the model representations. 

% \subsection{Background \& related work}

% Our work is an extension of recent theoretical and empirical results on embedding-based representations of generative models \citep{aranyak}, which itself is a continuation of a long-line of embedding-based investigations of the inputs and outputs of generative models \citep{mikolov2013efficient, reimers2019sentence, neelakantan2022text, 10098736}. 
% Of particular relevance is recent empirical work defining the data kernel \citep{duderstadt2023comparing} and investigations into its ability to track the dynamics of interacting models \citep{helm2024tracking}, in which the experiments demonstrate the ability to parlay the embeddings of a collection of inputs or outputs into useful vector representations of the generative models themselves in both white-box and black-box settings.
% The results herein further theoretically and empirically validate these findings in the context of statistical inference.

% Our work is also related to using embedding-based techniques for inference on complicated objects such as entire mouse connectomes \citep{WANG2020117274}, physiological data \citep{chen2022mental}, and classification distributions \citep{helm2021inducing}.
% In each of these settings, the authors define a pairwise distance matrix on the objects and apply multi-dimensionsal scaling to obtain vector representations of each of the objects.
% Once there is one vector representation per object, standard inference methods for the specific task can be used.
% The method we study herein follows this general formula, with the additional complication that the objects are random mappings.

% Lastly, our work is a part of the relatively new literature on inference on generative models. 
% For example, FlashHELM \citep{perlitz2023efficient} uses the score on a subset of a benchmark such as HELM \citep{liang2022holistic} to quickly identify where a model fits on a leaderboard. 
% More recent work proposes scaling laws to predict the performance of base models on a suite of benchmarks as a function of training FLOPs \citep{ruan2024observational}. 
% % \textcolor{red}{REF} uses the log-probabilities of input tokens to identify if a model has had access to sensitive data in its training process.
% The method proposed herein does not assume access to a scoring function nor does it assume access to a function of the model's weights at inference time.
% % Our approach can operate in a black-box setting and relies only the outputs and scores of a subset of previously scored models. 
% Perhaps most importantly, our setting and method can be applied to general model-level prediction scenarios.

\section{Introduction}
Generative models have recently met or surpassed human-level standards on benchmarks across a range of tasks \citep{nori2023capabilitiesgpt4medicalchallenge, katz2024gpt, dubey2024llama}. 
While these claims warrant skepticism and further robustness evaluation \citep{ness2024medfuzzexploringrobustnesslarge}, the impressive capabilities have created a competitive environment for training state-of-the-art models and have inspired the development of complementary methods to adapt models to particular use cases.
For example, quantizing the weights of a neural model enables a trade-off of model precision with vRAM and disk space \citep{gholami2022survey} while methods such as Low Rank Adaptation (LoRA) \citep{hu2021lora} and prompt-tuning \citep{lester2021power} enable compute- and data-efficient model adaptation. 
Other methods such as retrieval-augmented generation \citep{rag}, model merging \citep{model-merging}, constrained decoding \citep{constrained-decoding}, etc. \citep{dettmers2024qlora, graphrag}, have similarly contributed to the rapid development of a large population of diverse and accessible models.

Each model in the population has an accompanying set of covariates – scores on benchmarks, training mixture proportions, model safety scores, etc. – that are a function of the model, the training set, the architecture, the retrieval database, or derivatives thereof. 
For a given model, the covariate of interest may not be available to the user or it may be too expensive to calculate.
For example, it may not be known if a proprietary model has been trained on copyrighted data.
Or, it may be too resource intensive to directly evaluate the toxicity of every model.
Methods for predicting model-level
covariates are necessary in these settings, and others, to fully understand the behavior and properties
of a model.

In this paper we extend recent theoretical results for vector representations of black-box generative models \citep{aranyak} to model-level statistical inference settings.
Our results show that the embeddings of the responses from a collection of generative models can be used for consistent inference for a wide class of inference problems.
We demonstrate the effectiveness of the representations for three downstream inference tasks, including predicting the presence of sensitive information in a model's training mixture and predicting a model's safety. 
We include empirical investigations of performance sensitivity to hyperparameters required to generate the model-level vector representations. 

\noindent \textbf{Contribution.} Our contribution is the theoretical and empirical validation for using vector representations of black-box generative models for model-level inference tasks.
The representations that we study herein are based on the embeddings of their responses.
The theoretical results apply to any generic, well-behaved embedding function.
Given the representations, any vector-based method can be used for inference on the models.

\subsection{Background \& related work}

Our work is an extension of recent theoretical and empirical results on embedding-based representations of generative models \citep{aranyak}, which itself is a continuation of a long-line of embedding-based investigations of the inputs and outputs of generative models \citep{mikolov2013efficient, reimers2019sentence, neelakantan2022text, 10098736}. 
Of particular relevance is recent empirical work defining the data kernel \citep{duderstadt2023comparing} and investigations into its ability to track the dynamics of interacting models \citep{helm2024tracking}, in which the experiments demonstrate the ability to parlay the embeddings of a collection of inputs or outputs into useful vector representations of the generative models themselves in both white-box and black-box settings.
The results herein further theoretically and empirically validate these findings in the context of statistical inference.

Our work is also related to using embedding-based techniques for inference on complicated objects such as entire mouse connectomes \citep{WANG2020117274}, physiological data \citep{chen2022mental}, and classification distributions \citep{helm2021inducing}.
In each of these settings, the authors define a pairwise distance matrix on the objects and apply multi-dimensionsal scaling to obtain vector representations of each of the objects.
Once there is one vector representation per object, standard inference methods for the specific task can be used.
The method we study herein follows this general formula, with the additional complication that the objects are random mappings.

Lastly, our work is a part of the relatively new literature on inference on generative models. 
For example, FlashHELM \citep{perlitz2023efficient} uses the score on a subset of a benchmark such as HELM \citep{liang2022holistic} to quickly identify where a model fits on a leaderboard. 
More recent work proposes scaling laws to predict the performance of base models on a suite of benchmarks as a function of training FLOPs \citep{ruan2024observational}. 
% \textcolor{red}{REF} uses the log-probabilities of input tokens to identify if a model has had access to sensitive data in its training process.
The method proposed herein does not assume access to a scoring function nor does it assume access to a function of the model's weights at inference time.
% Our approach can operate in a black-box setting and relies only the outputs and scores of a subset of previously scored models. 
Perhaps most importantly, our setting and method can be applied to general model-level prediction.